\documentclass[12pt, a4paper]{article}

% --- Pacchetti Scientifici e Formattazione ---
\usepackage{amsmath, amssymb, amsfonts, amsthm}
\usepackage{booktabs, graphicx}
\usepackage[expansion=false]{microtype}
\usepackage[document]{ragged2e}
\usepackage{xcolor}

% --- Codifica e Font ---
\usepackage[utf8]{inputenc}
\usepackage[T1]{fontenc}
\usepackage[italian]{babel}
\usepackage{mathptmx}

% --- Gestione Margini e Layout ---
\usepackage[includeheadfoot]{geometry}
\geometry{
    a4paper,
    top=3cm,
    bottom=3cm,
    left=3.5cm,
    right=3.5cm,
}

% --- Gestione Interlinea e Figure ---
\usepackage{setspace}
\setstretch{1.3}
\usepackage[figuresright]{rotating}
\usepackage[pdftex, hidelinks]{hyperref}


\usepackage{subfiles}

% Document metadata
\title{Progettazione configuratore commerciale}
\author{Mattia Cavina}
\date{09 febbraio 2026}

\begin{document}
\maketitle
\newpage

% --- Table of Contents ---
\tableofcontents

% --- Abstract---
\documentclass[12pt, a4paper]{article} % Classe: article, report o book [8]

% --- Pacchetti Scientifici/Matematici ---
\usepackage{amsmath, amssymb, amsfonts, amsthm} % Matematica avanzata [2]
\usepackage{booktabs} % Per tabelle professionali

% --- Formattazione Testo e Altro ---
\usepackage[document]{ragged2e}
\usepackage{xcolor}    % Per usare i colori 

% --- Codifica e Font ---
\usepackage[utf8]{inputenc}
\usepackage[T1]{fontenc}
\usepackage[italian]{babel}

% Scelta del font (Times)
\usepackage{mathptmx} 

% --- Gestione Margini e Layout (Standard 65-70 caratteri) ---
\usepackage[includeheadfoot]{geometry}
\geometry{
    a4paper,
    top=3cm,
    bottom=3cm,
    left=3.5cm,   % Margine più ampio per la rilegatura
    right=3.5cm,  % Regolando questi margini si ottengono ~65-70 caratteri
}

% --- Gestione Righe (32-35 righe per pagina) ---
\usepackage{setspace}
\setstretch{1.3} % Interlinea personalizzata per circa 33-34 righe

% --- Gestione Figure e Tavole ---
\usepackage{graphicx}
\usepackage[figuresright]{rotating} % Per tavole in A3 o orizzontali

% --- PDF e Hyperlinks ---
\usepackage[pdftex, hidelinks]{hyperref}

% Document metadata
\title{Risultato delle interviste}
\author{Mattia Cavina}
\date{03 febbraio 2026}

\begin{document}
Il presente progetto di tesi riguarda la progettazione e implementazione di un nuovo configuratore commerciale per l'azienda Cepi~\cite{Cepi}. 
Attualmente l'azienda utilizza un programma sviluppato internamente basato su Access, il quale presenta significative limitazioni tecniche e funzionali dopo circa 30 anni di utilizzo. 
L'analisi si concentra sulla valutazione del sistema attuale, sulla raccolta sistematica dei requisiti attraverso interviste con gli utenti, e sulla progettazione di un nuovo database relazionale.
Il progetto è stato realizzato in collaborazione con il project manager esterno Claudio Rossi, che ha coordinato l'avanzamento dei lavori e ha contribuito all'implementazione della soluzione.
\end{document}

% --- Situazione As-Is ---
% \subfile{./DocumentazioneRaccolta/Analisi/As-Is_Access.tex}
\section{Situazione As-Is}
\subsection{Situazione Infrastrutturale Aziendale}
\subsection{Attuali software utilizzati}
\subsection{Organizzazione Aziendale e Flussi di lavoro}

% --- Interviste per la raccolta requisiti ---
% \subfile{./DocumentazioneRaccolta/Analisi/Raccolta_Requisiti.tex}
\section{Raccolta dei requisiti}
\subsection{Scelta dei candidati}
\subsection{Modalità di svolgimento raccolta informazioni}
\subsection{Riassunto punti emersi per reparto}

% --- Analisi dei requisiti emersi dalle interviste ---
%\subfile{./DocumentazioneRaccolta/Analisi/Analisi_Requisiti.tex}
\section{Analisi dei requisiti}
\subsection{Requisiti funzionali}
\subsection{Requisiti non funzionali}
\subsection{Requisiti futuri should-have}

\newpage
\section{Progettazione dell'architettura}
\subsection{Preambolo sulle scelte architetturali}
\subsection{Deployment}
\subsubsection{Diagramma dei Deployment}
\subsection{Componenti}
\subsubsection{Diagramma dei Componenti}
\subsection{Flussi di lavoro}
\subsubsection{Diagrammi di sequenza}

\newpage

\section{Scelte progettuali}
\subsection{Scelta modello di sviluppo}
\subsection{Suddivisione e contenuto dei rilasci}
\subsection{Timeline di progetto e successivi focus group}
\subsection{Scelta Framework ORM}
\subsection{Scelta Framework grafico} % MAUI .NET
\subsection{Scelta del pattern di progettazione} % MVVM

\newpage
% --- Progettazione del database ---
%\subfile{./DocumentazioneRaccolta/Database/Progettazione_database.tex}
\section{Progettazione del database}
\subsection{Esportazione Vecchio Db}
\subsection{Creazione regole di denominazione}
\subsection{Utilizzo di SSIS nel processo di migrazione}

\subsection{Ristrutturazione database}
\subsubsection{Eliminazione ridondanze e normalizzazione}
\subsubsection{Creazione tabelle di lookup}
\subsubsection{Aggiunta tabelle non presenti}
\subsubsection{Creazione tabelle di Back Up}
\subsubsection{Scelta di strutture Avanzate come Triggers e Stored Procedure}

\subsection{Schema logico finale}
\subsection{Popolamento database}
\subsection{Dati fissi}
\subsection{Dati fissi da vagliare}
\subsection{Dati provenienti da altri DB}
\subsubsection{Scelta dei dati}
\subsubsection{Query di sincronizzazione e tempistiche di aggiornamento}

\newpage
\section{Creazione Model e View}
\subsection{Mapping delle classi}
\subsection{Struttura del package}
\subsubsection{Classi di Test}
\subsubsection{Repository}
\subsubsection{Classi ed Interfacce esposte}

\subsection{Utilizzo di Figma per disegno interfacce}
\subsection{Creazione delle Factory per le View}
\subsection{Scelta delle classi di viewmodel}

\subsection{Utilizzo della Dependency Injection}
\subsubsection{Utilizzo di multi Threading per la gestione dei dati}
\subsection{Scelta gestione transazione e dei commit}
\subsection{Utilizzo della Cache in memoria per la gestione di dati fissi}

\newpage
\section{Modulo configurazione prodotto}
\subsection{Scopo del modulo}
\subsection{Focus group individuazione metodologia}
\subsection{Studio sulla guida alla configurazione}
\subsection{Query utilizzate per la ricerca}
\bibliographystyle{plain}
\bibliography{references}

\end{document}